\chapter{Control Flow \& Conditionals}

Programs need to make decisions based on changing data. In Python, \textbf{Control Flow} allows you to execute specific blocks of code depending on whether a given condition evaluates to \texttt{True} or \texttt{False}.

\section{Conditional Statements (\texttt{if}, \texttt{elif}, \texttt{else})}

Python uses indentation (4 spaces) to define code blocks instead of curly braces \texttt{\{\}}.

\begin{theorembox}{Syntax of Conditional Blocks}
\begin{lstlisting}
if condition_1:
    # Executed if condition_1 is True
elif condition_2:
    # Executed if condition_1 is False AND condition_2 is True
else:
    # Executed if all above conditions are False
\end{lstlisting}
\end{theorembox}

\begin{examplebox}{Categorizing Credit Scores}
\begin{lstlisting}
credit_score = 720

if credit_score >= 750:
    status = "Excellent"
elif credit_score >= 650:
    status = "Good"
elif credit_score >= 550:
    status = "Fair"
else:
    status = "Poor"

print(f"Credit Status: {status}")  # Output: Credit Status: Good
\end{lstlisting}
\end{examplebox}

\textbf{Data Science Relevance:} Decision Tree models (like Random Forests and XGBoost) work under the hood as a giant hierarchy of \texttt{if-elif-else} checks on feature columns!

\begin{videocard}{IITM BS Lecture: Introduction to the IF Statement}{https://www.youtube.com/watch?v=FTX5wF_3J9Q}
Official lecture explaining logical branching, indentation rules, and conditional execution.
\end{videocard}

\begin{videocard}{IITM BS Tutorial: If, Else, and Elif Conditions}{https://www.youtube.com/watch?v=-dBqiRCHbNw}
Hands-on tutorial with multiple practical examples of complex boolean expressions.
\end{videocard}

\section{Nested Conditionals and Short-Circuit Evaluation}

You can nest \texttt{if} blocks inside other \texttt{if} blocks. When evaluating compound boolean expressions using \texttt{and} and \texttt{or}, Python uses \textbf{Short-Circuit Evaluation}:
\begin{itemize}
    \item In \texttt{A and B}, if \texttt{A} is \texttt{False}, Python stops immediately and returns \texttt{False} without evaluating \texttt{B}.
    \item In \texttt{A or B}, if \texttt{A} is \texttt{True}, Python stops immediately and returns \texttt{True} without evaluating \texttt{B}.
\end{itemize}

\section{Importing Standard Libraries}

Python comes with a rich standard library of useful modules (like \texttt{math} and \texttt{random}). You can import them using different import syntaxes.

\begin{definitionbox}{Library Import Syntaxes}
\begin{lstlisting}
import math
print(math.sqrt(16))  # 4.0

import random as rd  # Using an alias
print(rd.randint(1, 6)) # Random integer from 1 to 6

from math import pi, log
print(log(10))
\end{lstlisting}
\end{definitionbox}

\begin{videocard}{IITM BS Lecture: Importing Libraries in Python}{https://www.youtube.com/watch?v=eW58_ky7oc8}
Official lecture comparing \texttt{import module}, \texttt{import module as alias}, and \texttt{from module import function}.
\end{videocard}

\newpage
\section{Practice Exercises}
\textit{Detailed step-by-step solutions are available in the Appendix.}

\begin{enumerate}
    \item \textbf{Leap Year Logic:} A year is a leap year if it is divisible by 4, EXCEPT if it is divisible by 100 (unless it is also divisible by 400). Write a Python script with an \texttt{if-else} statement to determine whether \texttt{year = 2024} is a leap year.

    \item \textbf{Short-Circuit Evaluation:} What will be printed by the following code snippet?
\begin{lstlisting}
x = 0
if x != 0 and (10 / x > 2):
    print("Valid")
else:
    print("Safe")
\end{lstlisting}
    Explain why Python does not crash with a \texttt{ZeroDivisionError}.

    \item \textbf{Ternary Operator (Conditional Expression):} In Python, you can write concise 1-line conditionals: \texttt{val = A if condition else B}. Convert the following block into a 1-line ternary expression:
\begin{lstlisting}
if score >= 50:
    result = "Pass"
else:
    result = "Fail"
\end{lstlisting}

    \item \textbf{Random Sampling:} Write a 3-line script that imports the \texttt{random} library and simulates rolling two 6-sided dice, printing their individual values and total sum.

    \item \textbf{Data Science Application (Anomaly Detection):} An IoT sensor sends temperature readings. If the temperature is between $18^\circ$C and $28^\circ$C (inclusive), it is \texttt{"Normal"}. If it is below $18^\circ$C, it is \texttt{"Warning: Low Temp"}. If it exceeds $28^\circ$C, it is \texttt{"Alert: Overheating"}. Write the \texttt{if-elif-else} block for \texttt{temp = 31.5}.
\end{enumerate}